\chapter{高等数学：一元函数微分学}

\section{导数定义、可导性与曲率}

\begin{problemblock}
\textbf{例1.} 已知函数 \(f(x)\) 在 \(x=1\) 处可导，且
\[
\lim_{x\to0}\frac{f(\cos x)-3f(1+\sin^2x)}{x^2}=2,
\]
求 \(f'(1)\)。
\end{problemblock}
\begin{solutionblock}
\analysis{分母是 \(x^2\)，先由极限有限推出常数项必须消失，再用 \(\cos x-1\) 与 \(\sin^2x\) 的二阶主部。}
当 \(x\to0\) 时，分子趋于 \(f(1)-3f(1)=-2f(1)\)。原极限有限，故 \(f(1)=0\)。
由 \(f\) 在 1 处可导，
\[
f(\cos x)=f'(1)(\cos x-1)+o(\cos x-1),\qquad
f(1+\sin^2x)=f'(1)\sin^2x+o(\sin^2x).
\]
又 \(\cos x-1\sim-\frac{x^2}{2}\)，\(\sin^2x\sim x^2\)，所以
\[
f(\cos x)-3f(1+\sin^2x)
=\left(-\frac12-3\right)f'(1)x^2+o(x^2)
=-\frac72f'(1)x^2+o(x^2).
\]
于是 \(-\frac72f'(1)=2\)，故
\[
f'(1)=-\frac47.
\]
\answer{\(\displaystyle f'(1)=-\frac47\)。}
\method{方法二（换元主部法）}
令 \(u=\cos x-1\sim-x^2/2\)，\(v=\sin^2x\sim x^2\)。因为 \(f(1)=0\)，原式等价于
\[
\frac{f'(1)u-3f'(1)v}{x^2}\to f'(1)\left(-\frac12-3\right),
\]
同样得到 \(f'(1)=-4/7\)。
\examnote{遇到 \(f(\cos x),f(1+\sin^2x)\)，先把内层都化为“1 加小量”，再用导数定义。}
\end{solutionblock}

\begin{problemblock}
\textbf{例2.} 设 \(y=f(x)\) 由方程
\[
x=\int_1^{y-x}\sin^2\left(\frac{\pi}{4}t\right)dt
\]
确定，求
\[
\lim_{n\to\infty}n\left(f\left(\frac1n\right)-1\right).
\]
\end{problemblock}
\begin{solutionblock}
\analysis{先求 \(f(0)\)，再对隐式方程求导。极限 \(n(f(1/n)-1)\) 本质就是 \(f'(0)\)。}
令 \(x=0\)，有
\[
0=\int_1^{y} \sin^2\left(\frac{\pi}{4}t\right)dt.
\]
被积函数非负且在 \(t=1\) 附近为正，故 \(y=1\)，即 \(f(0)=1\)。
设
\[
H(u)=\int_1^u\sin^2\left(\frac{\pi}{4}t\right)dt,
\]
则方程为 \(x=H(f(x)-x)\)。两边对 \(x\) 求导：
\[
1=H'(f(x)-x)(f'(x)-1).
\]
在 \(x=0\) 处，\(f(0)-0=1\)，且
\[
H'(1)=\sin^2\frac\pi4=\frac12.
\]
因此
\[
1=\frac12(f'(0)-1),\qquad f'(0)=3.
\]
于是
\[
\lim_{n\to\infty}n\left(f\left(\frac1n\right)-1\right)
=\lim_{h\to0^+}\frac{f(h)-f(0)}h=f'(0)=3.
\]
\answer{\(3\)。}
\examnote{隐函数给出的极限，常先识别成导数定义；本题的右端积分上限是 \(y-x\)，求导时别漏掉 \(-1\)。}
\end{solutionblock}

\begin{problemblock}
\textbf{例3.} 设函数 \(f(x)\) 在 \(x=0\) 处连续，且
\[
\lim_{x\to0}\frac{x f(x)-e^{2\sin x}+1}{\ln(1+x)+\ln(1-x)}=-3.
\]
证明 \(f(x)\) 在 \(x=0\) 处可导，并求 \(f'(0)\)。
\end{problemblock}
\begin{solutionblock}
\analysis{分母 \(\ln(1-x^2)\sim-x^2\)。分子中先由极限有限确定 \(f(0)\)，再比较一阶差商。}
由于
\[
e^{2\sin x}=1+2x+2x^2+o(x^2),
\]
分子为
\[
x f(x)-e^{2\sin x}+1=x(f(x)-2)-2x^2+o(x^2).
\]
分母
\[
\ln(1+x)+\ln(1-x)=\ln(1-x^2)=-x^2+o(x^2).
\]
若原极限有限，必须有 \(f(0)=2\)。于是
\[
\frac{x f(x)-e^{2\sin x}+1}{\ln(1+x)+\ln(1-x)}
= -\left(\frac{f(x)-2}{x}-2\right)+o(1).
\]
令其极限等于 \(-3\)，得
\[
-\lim_{x\to0}\frac{f(x)-2}{x}+2=-3,
\]
所以
\[
\lim_{x\to0}\frac{f(x)-f(0)}x=5.
\]
故 \(f\) 在 0 处可导，且 \(f'(0)=5\)。
\answer{\(f(0)=2\)，\(\displaystyle f'(0)=5\)。}
\method{方法二（待定系数法）}
设连续性给出 \(f(x)=A+Bx+o(x)\)。代入分子得 \((A-2)x+(B-2)x^2+o(x^2)\)。极限有限推出 \(A=2\)，再由 \(-(B-2)=-3\) 得 \(B=5\)。
\examnote{这类题的关键是“先消掉低阶项”。没有先推出 \(f(0)=2\)，直接谈 \(f'(0)\) 很容易漏逻辑。}
\end{solutionblock}

\begin{problemblock}
\textbf{例4.} 设
\[
f(x)=\begin{cases}
x^\alpha\sin x^\beta,&x\ne0,\\
0,&x=0,
\end{cases}
\]
其中 \(\alpha,\beta\) 为常数且 \(\beta<0\)。

(1) \(f(x)\) 在 \([-1,1]\) 上可导但 \(f'(x)\) 在 \([-1,1]\) 上无界，求 \(\alpha,\beta\) 满足的条件；

(2) \(f'(x)\) 在 \([-1,1]\) 上连续，求 \(\alpha,\beta\) 满足的条件。
\end{problemblock}
\begin{solutionblock}
\analysis{振荡来自 \(\sin x^\beta\)，因为 \(\beta<0\)，当 \(x\to0\) 时内层趋于无穷。先看可导性，再看导函数中的最坏项。}
由导数定义，
\[
f'(0)=\lim_{x\to0}\frac{x^\alpha\sin x^\beta}{x}
=\lim_{x\to0}x^{\alpha-1}\sin x^\beta.
\]
该极限存在且为 0 的条件是 \(\alpha>1\)。
当 \(x\ne0\) 时，
\[
f'(x)=\alpha x^{\alpha-1}\sin x^\beta+\beta x^{\alpha+\beta-1}\cos x^\beta.
\]
第一项在 \(\alpha>1\) 时趋于 0；第二项决定有界性和连续性。

(1) 要 \(f\) 可导，需 \(\alpha>1\)。要 \(f'\) 无界，需
\[
\alpha+\beta-1<0,
\]
即 \(\alpha+\beta<1\)。

(2) 要 \(f'\) 在 0 处连续，需第二项也趋于 0，即
\[
\alpha+\beta-1>0.
\]
由于 \(\beta<0\)，这自动推出 \(\alpha>1\)。
\answer{(1) \(\alpha>1,\ \alpha+\beta<1\)；(2) \(\alpha+\beta>1\)（且题设 \(\beta<0\)）。}
\examnote{振荡函数 \(x^p\sin x^\beta\) 中，\(p>0\) 控制连续，\(p>1\) 控制可导；求导后要额外检查链式法则带来的 \(x^{\beta-1}\)。}
\end{solutionblock}

\begin{problemblock}
\textbf{例5.} \(f(x)\) 在 \(x=0\) 处可导，且 \(f(x)\) 为偶函数，求 \(f'(0)\)。
\end{problemblock}
\begin{solutionblock}
\analysis{偶函数图像关于 \(y\) 轴对称，若在对称轴处可导，切线只能水平。}
由偶函数性质 \(f(-x)=f(x)\)。于是
\[
f'(0)=\lim_{x\to0}\frac{f(x)-f(0)}x,
\]
同时令 \(x\mapsto -x\)，得
\[
f'(0)=\lim_{x\to0}\frac{f(-x)-f(0)}{-x}
=-\lim_{x\to0}\frac{f(x)-f(0)}x=-f'(0).
\]
故 \(f'(0)=0\)。
\answer{\(0\)。}
\examnote{奇偶性与可导结合时，偶函数在 0 处导数为 0；奇函数若可导，则导函数为偶函数。}
\end{solutionblock}

\begin{problemblock}
\textbf{例6.} 设
\[
f(x)=(x-1)\sqrt{\frac{1-\sin x}{1+\cos x}}\cdot x^{x^2},
\]
求 \(f'(1)\)。
\end{problemblock}
\begin{solutionblock}
\analysis{函数含有因子 \((x-1)\)。求 \(x=1\) 处导数时，只需取其余因子在 1 处的值。}
记
\[
h(x)=\sqrt{\frac{1-\sin x}{1+\cos x}}\cdot x^{x^2}.
\]
则 \(f(x)=(x-1)h(x)\)，且 \(h\) 在 \(x=1\) 处连续。因此
\[
f'(1)=h(1).
\]
又 \(1^{1}=1\)，所以
\[
f'(1)=\sqrt{\frac{1-\sin 1}{1+\cos 1}}.
\]
\answer{\(\displaystyle \sqrt{\frac{1-\sin1}{1+\cos1}}\)。}
\examnote{乘积中有 \((x-a)\) 因子时，\([(x-a)h(x)]'_{x=a}=h(a)\)，这是考场快速题。}
\end{solutionblock}

\begin{problemblock}
\textbf{例7.} 证明：\(f(x)\) 在 \(x=x_0\) 处可导，\(g(x)\) 在 \(x=x_0\) 处连续但不可导，则
\[
F(x)=f(x)g(x)
\]
在 \(x=x_0\) 处可导当且仅当 \(f(x_0)=0\)。
\end{problemblock}
\begin{solutionblock}
\analysis{充分性用导数定义；必要性用反证：若 \(f(x_0)\ne0\)，则可由商 \(g=F/f\) 推出 \(g\) 可导。}
先证充分性。若 \(f(x_0)=0\)，则
\[
\frac{F(x)-F(x_0)}{x-x_0}
=\frac{f(x)g(x)}{x-x_0}
=\frac{f(x)-f(x_0)}{x-x_0}\,g(x).
\]
令 \(x\to x_0\)，由 \(f\) 可导、\(g\) 连续，得极限
\[
f'(x_0)g(x_0),
\]
所以 \(F\) 可导。

再证必要性。若 \(F\) 在 \(x_0\) 处可导且 \(f(x_0)\ne0\)，由于 \(f\) 可导从而连续，故 \(f\) 在 \(x_0\) 附近不为 0。于是
\[
g(x)=\frac{F(x)}{f(x)}.
\]
右端是两个在 \(x_0\) 处可导函数的商，分母非零，故 \(g\) 在 \(x_0\) 处可导，与题设矛盾。因此必须 \(f(x_0)=0\)。
\answer{\(F\) 在 \(x_0\) 处可导 \(\Longleftrightarrow f(x_0)=0\)。}
\examnote{“乘积可导但其中一个因子不可导”时，另一个因子必须在该点提供零因子来消掉不可导性。}
\end{solutionblock}

\begin{problemblock}
\textbf{例8.} 求
\[
F(x)=(x^2-x-2)|x^3-x|
\]
的不可导点。
\end{problemblock}
\begin{solutionblock}
\analysis{绝对值 \(|h(x)|\) 在 \(h\) 的单根处通常不可导；若外面还有零因子，可能把尖点消掉。}
令
\[
h(x)=x^3-x=x(x-1)(x+1),
\]
其零点为 \(-1,0,1\)，且都是单根。外因子
\[
x^2-x-2=(x-2)(x+1).
\]
在 \(x=0,1\) 处外因子不为 0，因此 \(|h(x)|\) 的尖点保留，\(F\) 不可导。
在 \(x=-1\) 处外因子也有因子 \((x+1)\)。局部看
\[
F(x)\sim C(x+1)|x+1|,
\]
而 \((x+1)|x+1|\) 在 \(-1\) 处导数存在且为 0，所以 \(-1\) 不是不可导点。
\answer{不可导点为 \(x=0,1\)。}
\method{方法二（分段法）}
也可按 \(x^3-x\) 的符号在区间 \((-
\infty,-1),(-1,0),(0,1),(1,+\infty)\) 上拆去绝对值，再只检查连接点 \(-1,0,1\)。
\examnote{绝对值题不要只看到 \(|\cdot|\) 的零点就全选；还要看外因子是否把尖点乘平。}
\end{solutionblock}

\begin{problemblock}
\textbf{例9.} 设 \(f(x)=3x^3+x^2|x|\)，则使得 \(f^{(n)}(0)\) 存在的最高阶数为（\quad）

A. 0 \qquad B. 1 \qquad C. 2 \qquad D. 3
\end{problemblock}
\begin{solutionblock}
\analysis{把含绝对值函数按 \(x\ge0\) 与 \(x<0\) 分段，比较左右高阶导数。}
当 \(x>0\) 时，\(f(x)=3x^3+x^3=4x^3\)；当 \(x<0\) 时，\(f(x)=3x^3-x^3=2x^3\)。
所以
\[
f'(0)=0,
\]
且左右一阶导数都趋于 0；二阶导数在 0 处也存在且为 0。
但三阶导数：右侧三阶导为 \(24\)，左侧三阶导为 \(12\)，左右不相等，故 \(f^{(3)}(0)\) 不存在。
\answer{C。最高阶数为 \(2\)。}
\examnote{形如 \(x^m|x|\) 的函数，在 0 附近等于左右系数不同的幂函数；最高可导阶通常卡在幂次处。}
\end{solutionblock}

\begin{problemblock}
\textbf{例10.} \(f(0)=0\)，则 \(f(x)\) 在 \(x=0\) 处可导的充分必要条件为（\quad）

A. \(\displaystyle\lim_{h\to0}\frac{f(1-\cos h)}{h^2}\) 存在；
B. \(\displaystyle\lim_{h\to0}\frac{f(1-e^h)}{h}\) 存在；

C. \(\displaystyle\lim_{h\to0}\frac{f(h-\sinh h)}{h^2}\) 存在；
D. \(\displaystyle\lim_{h\to\infty}\frac{f(2h)-f(h)}h\) 存在。
\end{problemblock}
\begin{solutionblock}
\analysis{可导等价于 \(\lim_{t\to0}f(t)/t\) 存在。看每个选项中的内层变量是否能等价覆盖 \(t\to0\) 的双侧邻域。}
因为 \(f(0)=0\)，\(f\) 在 0 处可导等价于
\[
\lim_{t\to0}\frac{f(t)}t
\]
存在。

B 中令 \(t=1-e^h\)，则 \(h\to0\) 时 \(t\to0\)，且
\[
t=1-e^h\sim-h.
\]
并且该变换在 0 附近一一覆盖 0 的双侧邻域。因此
\[
\frac{f(1-e^h)}h
=\frac{f(t)}t\cdot\frac{t}{h}\to -f'(0),
\]
存在性与 \(f'(0)\) 存在等价。
A 只对应 \(1-
\cos h\le0\)，主要检测单侧；C 中 \(h-\sinh h\sim-h^3/6\)，分母却是 \(h^2\)，不能等价为导数；D 与 \(x\to0\) 无关。
\answer{B。}
\method{方法二（反例排除）}
A 可取左右导数不同但左导数存在的函数排除；C 可取 \(f(t)=|t|^{2/3}\) 一类使选项极限存在但导数不存在；D 只约束无穷远处，不影响 0 点可导性。
\examnote{充分必要条件题要抓“变量替换是否双侧、同阶、可逆”。}
\end{solutionblock}

\begin{problemblock}
\textbf{例11.} 设函数 \(f(x)\) 连续，给出下列四个条件：
\[
\begin{array}{ll}
①\ \displaystyle\lim_{x\to0}\frac{|f(x)|-f(0)}x\text{ 存在},&
②\ \displaystyle\lim_{x\to0}\frac{f(x)-|f(0)|}x\text{ 存在},\\[2mm]
③\ \displaystyle\lim_{x\to0}\frac{|f(x)|}x\text{ 存在},&
④\ \displaystyle\lim_{x\to0}\frac{|f(x)|-|f(0)|}x\text{ 存在}.
\end{array}
\]
能得到 \(f(x)\) 在 \(x=0\) 处可导的条件个数是（\quad）

A. 1 \qquad B. 2 \qquad C. 3 \qquad D. 4
\end{problemblock}
\begin{solutionblock}
\analysis{分 \(f(0)>0,f(0)=0,f(0)<0\) 讨论。绝对值在非零点附近可去掉；在零点处，双侧极限会迫使差商为 0。}
①若极限存在，则必须 \(|f(0)|-f(0)=0\)，即 \(f(0)\ge0\)。若 \(f(0)>0\)，连续性保证 \(|f(x)|=f(x)\)，于是 \(f'(0)\) 存在；若 \(f(0)=0\)，则 \(|f(x)|/x\) 双侧极限存在，只能为 0，从而 \(|f(x)|=o(|x|)\)，故 \(f(x)/x\to0\)。①可推出可导。

②若 \(f(0)>0\)，它就是 \((f(x)-f(0))/x\)；若 \(f(0)=0\)，它就是 \(f(x)/x\)；若 \(f(0)<0\)，分子极限非零，不可能有有限极限。故②可推出可导。

③极限存在时，连续性迫使 \(f(0)=0\)。又 \(|f(x)|/x\) 双侧极限存在只能为 0，因此 \(f(x)/x\to0\)，可导。

④若 \(f(0)\ne0\)，绝对值在 0 附近等于 \(\pm f(x)\)，可推出 \(f'(0)\) 存在；若 \(f(0)=0\)，同③。故④也可推出可导。
\answer{D。四个条件均可推出 \(f\) 在 0 处可导。}
\examnote{绝对值差商的双侧极限非常强：当中心值为 0 时，左右符号相反会迫使极限只能为 0。}
\end{solutionblock}

\begin{problemblock}
\textbf{例12.} 设
\[
\lim_{x\to x_0^+}f'(x)=\lim_{x\to x_0^-}f'(x)=a,
\]
则（\quad）

A. \(f(x)\) 在 \(x=x_0\) 处必连续，但未必可导；
B. \(f(x)\) 在 \(x=x_0\) 处必可导且 \(f'(x_0)=a\)；

C. \(f(x)\) 在 \(x=x_0\) 处必有极限但未必连续；
D. 以上结论均不对。
\end{problemblock}
\begin{solutionblock}
\analysis{导函数在去心邻域有极限，可以控制 \(f(x)\) 的去心极限；但不能控制 \(f(x_0)\) 本身是否等于该极限。}
对 \(x\) 与 \(y\) 位于 \(x_0\) 同侧并趋于 \(x_0\)，由中值定理
\[
\frac{f(x)-f(y)}{x-y}=f'(\xi)	o a.
\]
可知 \(f(x)\) 在 \(x\to x_0^+\) 和 \(x\to x_0^-\) 时都有有限极限，且两侧极限一致，因此 \(\lim_{x\to x_0}f(x)\) 存在。
但 \(f(x_0)\) 可以另行定义。例如
\[
f(x)=x\ (x\ne0),\qquad f(0)=1,
\]
则 \(f'(x)=1\ (x\ne0)\)，两侧导数极限都为 1，但 \(f\) 在 0 处不连续，更不可导。
\answer{C。}
\examnote{去心邻域的导数信息不能自动修复中心点函数值；连续性仍要单独检查。}
\end{solutionblock}

\begin{problemblock}
\textbf{例13.} 设函数 \(f(x)\) 在 \((-
\infty,+\infty)\) 内有定义，在 \((-
\infty,0)\cup(0,+\infty)\) 内可导，且 \(\lim_{x\to0}f'(x)\) 存在，则 \(f(x)\) 在 \(x=0\) 处可导的一个充分条件是（\quad）

A. \(\displaystyle\lim_{x\to0}\frac{f(x)}x\) 存在；
B. \(\displaystyle\lim_{x\to0}\frac{f'(x)}x\) 存在；

C. \(f(x)\) 在 \(x=0\) 处连续；
D. \(\displaystyle\int_0^x f(t)dt\) 在 \(x=0\) 处可导。
\end{problemblock}
\begin{solutionblock}
\analysis{若中心点连续，再配合去心导数极限存在，就能用中值定理推出导数存在。其他选项都不能保证 \(f(0)\) 与去心极限相等。}
若 C 成立，即 \(f\) 在 0 处连续。对 \(x\ne0\)，由拉格朗日中值定理，存在 \(\xi\) 介于 0 与 \(x\) 之间，使
\[
\frac{f(x)-f(0)}x=f'(\xi).
\]
当 \(x\to0\) 时，\(\xi\to0\)，且 \(\lim_{t\to0}f'(t)\) 存在，所以差商极限存在，\(f\) 在 0 处可导。
A、B 都无法排除人为改变 \(f(0)\) 造成的不连续；D 只是积分平均意义上的条件，也不足以推出原函数在 0 处可导。
\answer{C。}
\examnote{“去心可导 + 导数极限存在”还差一个中心连续性，才能推出中心可导。}
\end{solutionblock}

\begin{problemblock}
\textbf{例14.} 设 \(f(x)\) 在 \(x=0\) 的邻域内二阶可导，曲线 \(y=f(x)\) 在 \((0,0)\) 处的曲率圆方程为
\[
(x-1)^2+(y-1)^2=2,
\]
求 \(f''(0)\)。
\end{problemblock}
\begin{solutionblock}
\analysis{曲率圆圆心为 \((1,1)\)，半径为 \(\sqrt2\)。半径方向与切线垂直，先由几何确定切线斜率，再用曲率公式。}
圆心 \((1,1)\) 到切点 \((0,0)\) 的半径向量为 \((-1,-1)\)，斜率为 1，因此切线斜率为 \(-1\)，即
\[
f'(0)=-1.
\]
曲率半径 \(R=\sqrt2\)，故曲率
\[
\kappa=\frac1R=\frac1{\sqrt2}.
\]
图形的主法向从 \((0,0)\) 指向 \((1,1)\)，与图像在 \(f'(0)=-1\) 时的正法向一致，因此取正号。由
\[
\kappa=\frac{f''(0)}{(1+[f'(0)]^2)^{3/2}},
\]
得
\[
f''(0)=\frac1{\sqrt2}(1+1)^{3/2}=2.
\]
\answer{\(2\)。}
\method{方法二（密切圆方向法）}
也可直接用有向曲率半径：中心位于单位法向 \((1,1)/\sqrt2\) 方向，故有向半径为 \(\sqrt2\)，有向曲率为 \(1/\sqrt2\)，代入同一公式得 \(f''(0)=2\)。
\examnote{曲率圆题要同时确定曲率大小和符号；只用半径只能得到 \(|f''(0)|\)。}
\end{solutionblock}

\begin{problemblock}
\textbf{例15.} 曲线
\[
e^x+6xy+y^2-1=0
\]
在点 \((0,0)\) 处的曲率半径为多少？
\end{problemblock}
\begin{solutionblock}
\analysis{在 \((0,0)\) 处 \(F_y=0,F_x\ne0\)，曲线切线为竖直方向，宜把 \(x\) 看作 \(y\) 的函数。}
将 \(x\) 写成 \(x=\varphi(y)\)。在 \((0,0)\) 附近设
\[
x=ay^2+o(y^2),
\]
因为切线竖直，所以没有一次项。代入
\[
e^x+6xy+y^2-1=0.
\]
由 \(e^x=1+x+o(x)\)，且 \(xy=O(y^3)\)，得
\[
1+ay^2+y^2-1+o(y^2)=0,
\]
故 \(a=-1\)。于是
\[
x=-y^2+o(y^2),\qquad \varphi''(0)=-2.
\]
对曲线 \(x=\varphi(y)\)，曲率
\[
\kappa=\frac{|\varphi''(0)|}{(1+[\varphi'(0)]^2)^{3/2}}=2.
\]
所以曲率半径
\[
R=\frac1\kappa=\frac12.
\]
\answer{\(\displaystyle \frac12\)。}
\method{方法二（隐函数求导）}
也可对 \(F(x,y)=0\) 视作 \(x=\varphi(y)\) 求导：\(F_x\varphi'+F_y=0\) 得 \(\varphi'(0)=0\)；再求导得 \(F_x\varphi''+F_{yy}=0\)，即 \(1\cdot\varphi''+2=0\)，故 \(\varphi''=-2\)。
\examnote{遇到竖直切线，换成 \(x=x(y)\) 计算曲率，能避免 \(y'\) 无穷大的麻烦。}
\end{solutionblock}

\begin{problemblock}
\textbf{例16.} 设函数 \(f(x),g(x)\) 在 \(x=0\) 处二阶导数连续，则
\[
\lim_{x\to0}\frac{f(x)-g(x)}{x^2}=0
\]
是两曲线 \(y=f(x),y=g(x)\) 在 \(x=0\) 对应的点处相切及曲率相等的（\quad）

A. 充分不必要条件 \qquad B. 充分必要条件

C. 必要不充分条件 \qquad D. 既不必要也不充分条件
\end{problemblock}
\begin{solutionblock}
\analysis{极限为 0 会推出函数值、一阶导数、二阶导数都相等，因此足够；但曲率是非负量，只要求二阶导数绝对意义相配，不一定要求二阶导数相等。}
由 Taylor 展开，
\[
f(x)-g(x)=[f(0)-g(0)]+[f'(0)-g'(0)]x+\frac{1}{2}[f''(0)-g''(0)]x^2+o(x^2).
\]
若
\[
\frac{f(x)-g(x)}{x^2}	o0,
\]
则必须有
\[
f(0)=g(0),\quad f'(0)=g'(0),\quad f''(0)=g''(0).
\]
这显然推出两曲线在对应点处相切，且曲率
\[
\kappa=\frac{|y''|}{(1+(y')^2)^{3/2}}
\]
相等。因此该条件充分。
但它不必要。例如 \(f(x)=x^2,\ g(x)=-x^2\)，两曲线在 0 处有同一点、同切线，曲率都为 2；但
\[
\frac{f(x)-g(x)}{x^2}=2\not\to0.
\]
\answer{A。充分不必要条件。}
\examnote{“曲率相等”通常指曲率大小相等，不含凹向符号；而二阶 Taylor 主项相等比曲率相等更强。}
\end{solutionblock}

\begin{problemblock}
\textbf{例17.} 设曲线 \(L\) 由参数方程
\[
\begin{cases}
x=\frac13t^3+t^2+t,\\
y=t^4,
\end{cases}\qquad (t\ge -2)
\]
确定，则曲线上的拐点是哪里？
\end{problemblock}
\begin{solutionblock}
\analysis{参数曲线的凹向由 \(d^2y/dx^2\) 的符号判断，注意 \(x'(t)=0\) 处可能是竖直切线，也可能发生凹向改变。}
计算
\[
x'(t)=t^2+2t+1=(t+1)^2,
\quad x''(t)=2(t+1),
\]
\[
y'(t)=4t^3,
\quad y''(t)=12t^2.
\]
当 \(t\ne-1\) 时，
\[
\frac{d^2y}{dx^2}
=\frac{x'y''-y'x''}{(x')^3}
=\frac{4t^2(t+1)(t+3)}{(t+1)^6}
=\frac{4t^2(t+3)}{(t+1)^5}.
\]
在区间 \(t\ge-2\) 内，\(t=-1\) 两侧该式符号由负变正；\(t=0\) 虽使二阶导为 0，但符号不变。因此拐点对应 \(t=-1\)。
代回得
\[
x=\frac13(-1)^3+(-1)^2-1=-\frac13,
\qquad y=(-1)^4=1.
\]
\answer{\(\displaystyle\left(-\frac13,1\right)\)。}
\examnote{参数方程题不要只解 \(d^2y/dx^2=0\)，还要检查 \(x'(t)=0\) 的竖直切线点和符号变化。}
\end{solutionblock}

\begin{problemblock}
\textbf{例18.} 设
\[
e^y=3+7x-6x^2,
\]
求 \(y^{(n)}\)。
\end{problemblock}
\begin{solutionblock}
\analysis{先对右端二次式因式分解，把 \(y\) 化成对数之和，再使用 \(\ln(1+ax)\) 的高阶导数公式。}
由
\[
3+7x-6x^2=3(1+3x)\left(1-\frac23x\right),
\]
得
\[
y=\ln3+
\ln(1+3x)+\ln\left(1-\frac23x\right).
\]
对 \(n\ge1\)，
\[
\frac{d^n}{dx^n}\ln(1+ax)=(-1)^{n-1}(n-1)!\frac{a^n}{(1+ax)^n},
\]
\[
\frac{d^n}{dx^n}\ln(1-bx)=-(n-1)!\frac{b^n}{(1-bx)^n}.
\]
因此
\[
y^{(n)}=(n-1)!\left[(-1)^{n-1}\frac{3^n}{(1+3x)^n}
-\frac{(2/3)^n}{\left(1-\frac{2}{3}x\right)^n}\right].
\]
若题目要求在 \(x=0\) 处的值，则
\[
y^{(n)}(0)=(n-1)!\left[(-1)^{n-1}3^n-\left(\frac23\right)^n\right].
\]
\answer{\(\displaystyle y^{(n)}=(n-1)!\left[(-1)^{n-1}\frac{3^n}{(1+3x)^n}-\frac{(2/3)^n}{(1-2x/3)^n}\right]\)。}
\examnote{高阶导数题先拆成基本函数；对数型最常用的是 \(\ln(1+ax)\) 的公式。}
\end{solutionblock}

\begin{problemblock}
\textbf{例19.} 函数 \(y=y(x)\) 由方程
\[
|x|y+e^y=e
\]
确定，则下列关于点 \(x=0\) 的说法正确的是（\quad）

A. 不可导点，\(y(0)\) 是极大值；
B. 不可导点，\(y(0)\) 不是极值；

C. 可导点，\(y(0)\) 不是极值；
D. 可导点，\(y(0)\) 是极大值。
\end{problemblock}
\begin{solutionblock}
\analysis{绝对值使左右导数不同。先求 \(y(0)\)，再分别在 \(x>0\)、\(x<0\) 两侧求导。}
令 \(x=0\)，得 \(e^{y(0)}=e\)，所以 \(y(0)=1\)。
当 \(x>0\) 时，方程为
\[
xy+e^y=e.
\]
求导得
\[
y+xy'+e^y y'=0,
\]
令 \(x\to0^+\)，得
\[
1+e y'_+(0)=0,
\qquad y'_+(0)=-\frac1e.
\]
当 \(x<0\) 时，方程为
\[
-xy+e^y=e.
\]
求导得
\[
-y-xy'+e^y y'=0,
\]
令 \(x\to0^-\)，得
\[
-1+e y'_-(0)=0,
\qquad y'_-(0)=\frac1e.
\]
左右导数不等，故不可导。又左侧靠近 0 时 \(y\) 随 \(x\) 增大而增大，右侧靠近 0 时 \(y\) 随 \(x\) 增大而减小，所以 \(y(0)=1\) 是极大值。
\answer{A。}
\examnote{含 \(|x|\) 的隐函数题必须左右分开求导，不能直接把 \(|x|\) 当普通可导函数处理。}
\end{solutionblock}

\begin{problemblock}
\textbf{例20.} \(f(x)\) 在 \(x=0\) 的邻域内有定义，去心邻域内可导，则下列说法中正确的是（\quad）

A. 若 \(\displaystyle\lim_{x\to0}\frac{f(x)}x=m\)，则 \(\displaystyle\lim_{x\to0}f(x)=0\) 且 \(f'(0)=m\)；

B. 若 \(\displaystyle\lim_{x\to0}f(x)=0\) 且 \(f'(0)=m\)，则 \(\displaystyle\lim_{x\to0}\frac{f(x)}x=m\)；

C. 若 \(\displaystyle\lim_{x\to0}\frac{f(x)}{x^2}=m\)，则 \(\displaystyle\lim_{x\to0}\frac{f'(x)}x=2m\)；

D. 若 \(\displaystyle\lim_{x\to0}\frac{f'(x)}x=2m\)，则 \(\displaystyle\lim_{x\to0}\frac{f(x)}{x^2}=m\)。
\end{problemblock}
\begin{solutionblock}
\analysis{本题考查“函数极限、导数、洛必达逆命题”的逻辑强弱。只要中心函数值未被控制，很多看似自然的命题都不成立。}
B 正确。若 \(f'(0)=m\)，则 \(f\) 在 0 处连续，所以 \(f(0)=\lim_{x\to0}f(x)=0\)。由导数定义
\[
\lim_{x\to0}\frac{f(x)}x
=\lim_{x\to0}\frac{f(x)-f(0)}{x-0}=f'(0)=m.
\]
A 错在未给 \(f(0)=0\)。例如令 \(f(x)=mx\ (x\ne0), f(0)=1\)，则 \(f(x)/x\to m\)，但 \(f\) 在 0 处不可导。
C 是洛必达逆命题，错误；可构造 \(f(x)=mx^2+x^3\sin(1/x^2)\) 破坏 \(f'(x)/x\) 的极限。
D 同样缺少中心值控制。即使 \(f'(x)/x\to2m\)，也可能有非零常数项，使 \(f(x)/x^2\) 不存在。
\answer{B。}
\examnote{洛必达法则不能随便反用；涉及 \(f'(0)\) 时还必须核对 \(f(0)\)。}
\end{solutionblock}

\begin{problemblock}
\textbf{例21.} 设函数 \(f(x)\) 二阶可导，则（\quad）

A. 若 \(f(x)\) 在 \(x=x_0\) 的某邻域内严格单调递增，则 \(f'(x_0)>0\)；

B. 若 \(f''(x_0)>0\)，则 \(f(x)\) 在 \(x=x_0\) 的某邻域内是凹的；

C. 若 \((x_0,f(x_0))\) 是 \(f(x)\) 的拐点，则 \(x=x_0\) 是 \(f'(x)\) 的极值点；

D. 若 \(x=x_0\) 是 \(f'(x)\) 的极值点，则 \((x_0,f(x_0))\) 是 \(f(x)\) 的拐点。
\end{problemblock}
\begin{solutionblock}
\analysis{按考研常用定义，拐点对应凹凸性改变，也就是 \(f'\) 的单调性在该点两侧发生由增到减或由减到增的变化。}
A 错。例如 \(f(x)=x^3\) 在 0 附近严格递增，但 \(f'(0)=0\)。
B 错。仅有某一点 \(f''(x_0)>0\) 不足以推出某邻域内 \(f''>0\)，除非额外给出 \(f''\) 连续。
C 正确。拐点意味着曲线凹凸性改变；对可导函数来说，就是 \(f'\) 在该点两侧的单调性发生改变，因此 \(f'\) 在 \(x_0\) 处取得局部极值。
D 不作为必然结论。\(f'\) 在一点有极值若缺少凹凸性改变的严格条件，不能直接推出该点一定为拐点。
\answer{C。}
\examnote{本题按教材型“凹凸性与导函数单调性”的对应关系判断；点值条件如 \(f''(x_0)>0\) 不能代替邻域条件。}
\end{solutionblock}

\begin{problemblock}
\textbf{例22.} 曲线
\[
y=(x-1)(x-2)^2(x-3)^3(x-4)^4
\]
的拐点是（\quad）

A. \((1,0)\) \qquad B. \((2,0)\) \qquad C. \((3,0)\) \qquad D. \((4,0)\)
\end{problemblock}
\begin{solutionblock}
\analysis{候选点都是零点，但拐点要看二阶导或凹凸性是否变号，不能只看函数是否过轴。}
设
\[
y=(x-a)^m\phi(x),\qquad \phi(a)\ne0.
\]
若 \(m\ge3\) 为奇数，则在 \(x=a\) 附近二阶导的主部含 \((x-a)^{m-2}\)，会随 \(x-a\) 变号，通常产生拐点。
本题各根重数分别为
\[
1,2,3,4.
\]
只有 \(x=3\) 是奇重数且重数至少为 3 的零点。直接看局部主部：
\[
y\sim C(x-3)^3\quad(C\ne0),
\]
故二阶导主部 \(6C(x-3)\) 变号，\((3,0)\) 为拐点。其他点不满足凹凸性变号。
\answer{C。\((3,0)\)。}
\examnote{多项式零点判拐点时，三重根、五重根等奇重高阶根最敏感；单根只是穿轴，不一定是拐点。}
\end{solutionblock}

\begin{problemblock}
\textbf{例23.} 设 \(y=f(x)\) 在 \((-
\infty,\infty)\) 上连续，且其导函数 \(f'(x)\) 的图像如原图所示，其中 \(x=0\) 和 \(x=x_5\) 是 \(f'(x)\) 的铅直渐近线。求 \(y=f(x)\) 极值点的个数与拐点的个数。
\end{problemblock}
\begin{solutionblock}
\analysis{极值看 \(f'\) 的符号变化；拐点看 \(f'\) 的增减变化，即 \(f''\) 的符号变化。铅直渐近线也要检查两侧符号或增减是否发生改变。}
由图像可读出：\(f'\) 在 \(x_1\) 处由正变负，\(f\) 有极大值；在 \(x_3\) 处由负变正，\(f\) 有极小值；在 \(x=0\) 左侧 \(f'>0\)、右侧 \(f'<0\)，且 \(f\) 连续，故 \(x=0\) 也是极大值点；在 \(x_4\) 处 \(f'\) 由负变正，故有极小值。\(x_6\) 处 \(f'=0\) 但两侧符号不变，不是极值点。
所以极值点共
\[
4\text{ 个}.
\]
拐点看 \(f'\) 的增减：左侧抛物线在 \(x_2\) 处由减变增，对应一个拐点；在 \(x=x_5\) 处，\(f'\) 左侧增至 \(+\infty\)，右侧从 \(+\infty\) 减小，凹凸性改变；在 \(x_6\) 处 \(f'\) 由减变增，又对应一个拐点。\(x=0\) 两侧 \(f'\) 都表现为递增，不产生凹凸性改变。
故拐点共
\[
3\text{ 个}.
\]
\answer{极值点 \(4\) 个，拐点 \(3\) 个。}
\examnote{看 \(f'\) 图像：过零且变号对应 \(f\) 极值；\(f'\) 的极值或增减性改变对应 \(f\) 拐点。}
\end{solutionblock}

\begin{problemblock}
\textbf{例24.} 函数 \(y=f(x)\) 在 \((-
\infty,+\infty)\) 内连续，其二阶导函数的图像如原图所示，则 \(y=f(x)\) 的拐点个数是（\quad）

A. 1 \qquad B. 2 \qquad C. 3 \qquad D. 4
\end{problemblock}
\begin{solutionblock}
\analysis{直接看 \(f''\) 的符号变化。\(f''=0\) 但不变号的点不是拐点；\(f''\) 不存在但两侧变号的点也可能是拐点。}
由图可见：在 \(x_1\) 附近，\(f''\) 从负变正，产生一个拐点；在 \(x_3\) 处，\(f''=0\) 但左右同为正，不产生拐点；在 \(x=0\) 附近，\(f''\) 左侧为正、右侧为负，凹凸性改变，产生一个拐点；在 \(x_4\) 处，\(f''\) 从负变正，产生一个拐点。
所以拐点共有
\[
3\text{ 个}.
\]
\answer{C。}
\examnote{拐点判定只认凹凸性变号；单纯 \(f''=0\) 或图像碰到坐标轴但不穿过，都不能算。}
\end{solutionblock}
